\section{Prescribed middle factors of polynomial coefficients}

The algebraic statements in this section describe the ordered
products used in \cite[Proposition~5.7]{Arad2016RigorousRG}.
The middle operators are specified before the coefficients on the exterior
factor. All products retain their order.
An interval $\{0,\ldots,h-1\}$ is empty when $h=0$.

For a matrix $M$ indexed by a product $m\times n$, use the right operator
blocks and left coefficient slices
\begin{align}
    M^{\mathrm R}_{ab}(x,y) & =M((a,x),(b,y)),\notag \\
    M^{\mathrm L}_{xy}(a,b) & =M((a,x),(b,y)).
    \label{eq:two_marker_operator_blocks}
\end{align}
Here $a,b\in m$ and $x,y\in n$; the right blocks are indexed by $n$
and the left slices by $m$.

\begin{lemma}[Prescribed factors from operator blocks]
    \label{thm:two_marker_prescribed_span}
    \lean{Matrix.exists_left_factors_of_operatorBlock_mem_span,Matrix.exists_right_factors_of_left_blocks_mem_span}
    \leanok
    Let $\mathcal I$ be finite, let $m,n$ be arbitrary sets, and let
    $C_\alpha$ be complex matrices indexed by $n$. Suppose that
    $M^{\mathrm R}_{ab}\in\spn_{\C}\{C_\alpha:\alpha\in\mathcal I\}$ for every $a,b\in m$.
    There are matrices $L_\alpha$ indexed by $m$ such that
    \begin{align}
        M=\sum_{\alpha\in\mathcal I}L_\alpha\otimes C_\alpha.
        \label{eq:two_marker_prescribed_span}
    \end{align}
    For the opposite orientation, let $C_\alpha\in\C^{m\times m}$ for
    $\alpha\in\mathcal I$. If
    $M^{\mathrm L}_{xy}\in\spn_{\C}\{C_\alpha:\alpha\in\mathcal I\}$
    for every $x,y\in n$, then
    $M=\sum_{\alpha\in\mathcal I}C_\alpha\otimes R_\alpha$ for suitable
    matrices $R_\alpha$ indexed by $n$.
\end{lemma}
\begin{proof}\leanok
    For each $a,b$, choose coefficients $\lambda_{ab,\alpha}$ with
    $M^{\mathrm R}_{ab}=\sum_\alpha\lambda_{ab,\alpha}C_\alpha$, and set
    $(L_\alpha)_{ab}=\lambda_{ab,\alpha}$. Comparing matrix entries gives
    (\ref{eq:two_marker_prescribed_span}). The opposite orientation follows by
    the same argument with the two pairs of indices exchanged.
\end{proof}

\begin{definition}[Crossing words and ordinary run lengths]
    \label{def:two_marker_middle_index}
    \lean{TNLean.MPS.AreaLaw.TwoMarkerMiddleIndex}
    \leanok
    For nonnegative integers $r,i,j,u,v$, let $\mathcal I_{r,i,j}^{u,v}$ consist
    of pairs $(w,g)$ with
    \begin{align}
        w & :\{0,\ldots,i+j-1\}\longrightarrow
        \{0,\ldots,u-1\}\sqcup\{0,\ldots,v-1\},\notag       \\
        g & :\{0,\ldots,i+j\}\longrightarrow\{0,\ldots,r\}.
        \label{eq:two_marker_index_data}
    \end{align}
    The word has $i$ letters in the first summand and $j$ in the second, with
    $\sum_{h=0}^{i+j}g_h+i+j\le r$.
\end{definition}

\begin{definition}[Ordered middle product]
    \label{def:two_marker_middle_product}
    \lean{TNLean.MPS.AreaLaw.twoMarkerMiddleFactor}
    \leanok
    \uses{def:two_marker_middle_index}
    Let $A$ be an element of a monoid, and let $U_0,\ldots,U_{u-1}$ and
    $V_0,\ldots,V_{v-1}$ belong to the same monoid. For
    $(w,g)\in\mathcal I_{r,i,j}^{u,v}$, put $T(\operatorname{inl}s)=U_s$,
    $T(\operatorname{inr}t)=V_t$, and define
    \begin{align}
        \Phi_{w,g}(A,U,V)
        =\left(\prod_{h=0}^{i+j-1} A^{g_h}T(w_h)\right)A^{g_{i+j}}.
        \label{eq:two_marker_middle_product}
    \end{align}
    The product over $h$ is in increasing order; an empty product is the
    identity.
\end{definition}

\begin{definition}[Two distinguished frequencies]
    \label{def:two_marker_coefficient}
    \lean{TNLean.MPS.AreaLaw.twoMarkerCoefficient}
    \leanok
    Let $D,U,V$ belong to a semiring. With two central formal variables $x,y$,
    define
    \begin{align}
        C_{r,i,j}(D,U,V)=[x^i y^j](D+yV+xU)^r.
        \label{eq:two_marker_coefficient}
    \end{align}
    The variable $x$ counts $U$ and $y$ counts $V$; the coefficients $D,U,V$
    need not commute.
\end{definition}

\begin{lemma}[Operator blocks of sums and tensor products]
    \label{thm:two_marker_block_identities}
    \lean{TNLean.MPS.AreaLaw.operatorBlock_sum,TNLean.MPS.AreaLaw.operatorBlock_kronecker_eq_smul}
    \leanok
    For a finite family of matrices $M_\alpha$ indexed by $m\times n$, and
    matrices $E$ and $W$ indexed by $m$ and $n$, respectively,
    \begin{align}
        \left(\sum_\alpha M_\alpha\right)^{\mathrm R}_{ab}
                                      & =\sum_\alpha(M_\alpha)^{\mathrm R}_{ab},\notag \\
        (E\otimes W)^{\mathrm R}_{ab} & =E_{ab}W.
        \label{eq:two_marker_block_identities}
    \end{align}
    No finiteness assumption on $m$ or $n$ is required.
\end{lemma}
\begin{proof}\leanok
    Both identities follow by comparing the entry indexed by $x,y\in n$.
\end{proof}

\begin{lemma}[A common prescribed middle family through a degree bound]
    \label{thm:two_marker_common_middle}
    \lean{TNLean.MPS.AreaLaw.exists_sum_twoMarkerCoefficient_decomposition}
    \leanok
    \uses{def:two_marker_middle_index,def:two_marker_middle_product,def:two_marker_coefficient}
    Let $m,n$ be finite sets. Choose complex matrices $E,E_s^U,E_t^V$ indexed
    by $m$ and $A,U_s^M,V_t^M$ indexed by $n$, with $0\le s<u$ and $0\le t<v$.
    Put
    \begin{align}
        D & =E\otimes I+I\otimes A,\notag              \\
        U & =\sum_{s=0}^{u-1}E_s^U\otimes U_s^M,\notag \\
        V & =\sum_{t=0}^{v-1}E_t^V\otimes V_t^M.
        \label{eq:two_marker_split_letters}
    \end{align}
    For all nonnegative integers $k,i,j$ and all complex numbers
    $c_0,\ldots,c_k$, there are matrices $B_{w,g}$ indexed by $m$ such that
    \begin{align}
        \sum_{r=0}^k c_r C_{r,i,j}(D,U,V)
        =\sum_{(w,g)\in\mathcal I_{k,i,j}^{u,v}}
        B_{w,g}\otimes\Phi_{w,g}(A,U^M,V^M).
        \label{eq:two_marker_common_middle}
    \end{align}
    The displayed middle family depends only on $A,U^M,V^M,k,i,j$.
\end{lemma}
\begin{proof}\leanok
    \uses{thm:two_marker_prescribed_span,thm:two_marker_block_identities,def:two_marker_middle_index,def:two_marker_middle_product,def:two_marker_coefficient}
    Induction on $r$ shows that every right operator block of
    $C_{r,i,j}(D,U,V)$ belongs to the span of the products indexed by
    $\mathcal I_{r,i,j}^{u,v}$. At degree zero, the coefficient is $I$ for
    $i=j=0$ and zero otherwise. The positive-frequency recurrence is
    \begin{align}
        C_{r+1,i+1,j+1}
        =C_{r,i+1,j+1}D+C_{r,i+1,j}V+C_{r,i,j+1}U.
        \label{eq:two_marker_right_recurrence}
    \end{align}
    When one frequency vanishes, the corresponding last-letter term is absent.
    Right multiplication by $E\otimes I$ preserves a middle word and enlarges
    its degree bound. Right multiplication by $I\otimes A$ increases its final
    ordinary run. The $U$ and $V$ terms append the corresponding crossing
    letter, retaining the order in (\ref{eq:two_marker_middle_product}).
    Lemma~\ref{thm:two_marker_prescribed_span} and
    (\ref{eq:two_marker_block_identities}) justify these span operations.

    If $r\le k$, the same crossing letters and run lengths define an element
    of $\mathcal I_{k,i,j}^{u,v}$ with the same ordered product. Thus all the
    right operator blocks in the left side of
    (\ref{eq:two_marker_common_middle}) belong to the span of the one displayed
    middle family. Apply Lemma~\ref{thm:two_marker_prescribed_span} once more
    to obtain the matrices $B_{w,g}$.
\end{proof}

\begin{lemma}[Number of prescribed middle factors]
    \label{thm:two_marker_middle_cardinality}
    \lean{TNLean.MPS.AreaLaw.card_twoMarkerMiddleIndex_square_le}
    \leanok
    \uses{def:two_marker_middle_index}
    For all nonnegative integers $r,i,j,q$,
    \begin{align}
        |\mathcal I_{r,i,j}^{q^2,q^2}|
        \le 2^{i+j}q^{2(i+j)}(r+1)^{i+j+1}.
        \label{eq:two_marker_middle_cardinality}
    \end{align}
\end{lemma}
\begin{proof}\leanok
    \uses{def:two_marker_middle_index}
    There are at most $(2q^2)^{i+j}$ choices for the decorated word and
    $(r+1)^{i+j+1}$ choices for its run lengths. Discarding the frequency and
    total-length restrictions gives their product as an upper bound.
\end{proof}
